\documentclass[11pt]{article}
\usepackage{mathblatt}
\begin{document}
% Kopfzeile Seite 1 ohne Kennung; die Rückseite (Lösungen, für den Lehrer) bekommt sie nach \begleitteil
\blattkopf{Mathematik}{Basisaufgaben}
\einheitenkopf[][Zettel 13]{Basisaufgaben · Zettel 13}
\zweigzeile{je Aufgabe 1 Punkt · Lösungen auf der Rückseite}
\par\vspace{2pt}\noindent Name \underline{\hspace{7cm}}\hfill Datum \underline{\hspace{3cm}}\hfill Punkte \underline{\hspace{1.2cm}} / 18\par\vspace{6pt}
% \small: volle Seite, kurze Aufgaben zweispaltig (v0.6)
\small

\noindent\begin{minipage}[t]{0.485\linewidth}\raggedright
% 1: Mitte zweier Zahlen bestimmen – brueche-dezimalzahlen-basis-k3-v2 (Original 2020-OS-B1f, ausgedacht: keine Marke)
\begin{aufgabe}{Welche Zahl liegt genau in der Mitte zwischen $-2{,}1$ und $-2$? \leerfeld}
\end{aufgabe}
\end{minipage}\hfill\begin{minipage}[t]{0.485\linewidth}\raggedright
% 2: Termwert berechnen – rationale-zahlen-basis-k1-v9 (Original 2026-FOR-B1g, ausgedacht: keine Marke)
\begin{aufgabe}{Berechne den Wert von $-3 \cdot (x + 2)$ für $x = -5$. \leerfeld}
\end{aufgabe}
\end{minipage}\par

\noindent\begin{minipage}[t]{0.485\linewidth}\raggedright
% 3: Zahl zu Bedingung angeben – rationale-zahlen-basis-k3-v3 (Original 2015-OS-B1b, ausgedacht: keine Marke)
\begin{aufgabe}{Gib eine Zahl an, die zwischen $-2$ und $-1$ liegt. \leerfeld}
\end{aufgabe}
\end{minipage}\hfill\begin{minipage}[t]{0.485\linewidth}\raggedright
% 4: Zehnerpotenzschreibweise umwandeln – potenzen-wurzeln-basis-k3-v6 (Original 2025-OS-B1d, ausgedacht: keine Marke)
\begin{aufgabe}{$0{,}08 = 8 \cdot 10^{\square}$ – welche Hochzahl? \leerfeld}
\end{aufgabe}
\end{minipage}\par

\noindent\begin{minipage}[t]{0.485\linewidth}\raggedright
% 5: Prozentsatz berechnen – prozentrechnung-basis-k3-v2 (Original 2014-OS-B1e, ausgedacht: keine Marke)
\begin{aufgabe}{Von $400$ Schülern kommen $36$ mit dem Rad. Wie viel Prozent sind das? \leerfeld[\%]}
\end{aufgabe}
\end{minipage}\hfill\begin{minipage}[t]{0.485\linewidth}\raggedright
% 6: Prozentwert berechnen – prozentrechnung-basis-k4-v9 (Original 2026-FOR-B1a, ausgedacht: keine Marke)
\begin{aufgabe}{$60\,\%$ von $25\,€$ \leerfeld[€]}
\end{aufgabe}
\end{minipage}\par

\noindent\begin{minipage}[t]{0.485\linewidth}\raggedright
% 7: Größen vergleichen – einheiten-basis-k2-v2 (Original 2026-FOR-B1d, ausgedacht: keine Marke)
\begin{aufgabe}{Vergleiche $1{,}5$ kg und $1500$ g. Setze das richtige Zeichen ein: $<$, $=$ oder $>$. $1{,}5$ kg \leerfeld $1500$ g}
\end{aufgabe}
\end{minipage}\hfill\begin{minipage}[t]{0.485\linewidth}\raggedright
% 8: Portionen aus Gesamtmenge berechnen – einheiten-basis-k3-v1 (Original 2022-OS-B1h, ausgedacht: keine Marke)
\begin{aufgabe}{In einer Kanne sind $2$ l Saft. Ein Becher fasst $250$ ml. Wie viele Becher kann man füllen? \leerfeld Becher}
\end{aufgabe}
\end{minipage}\par

\noindent\begin{minipage}[t]{0.485\linewidth}\raggedright
% 9: Proportionale Zuordnung Dreisatz – zuordnungen-basis-k2-v2 (Original 2023-OS-B1a, ausgedacht: keine Marke)
\begin{aufgabe}{$5$ l Benzin kosten $9$ €. Wie viel kosten $8$ l? \leerfeld[€]}
\end{aufgabe}
\end{minipage}\hfill\begin{minipage}[t]{0.485\linewidth}\raggedright
% 10: Lösung durch Einsetzen prüfen – quadratische-gleichungen-basis-k1-v6 (Original 2025-OS-B1h, ausgedacht: keine Marke)
\begin{aufgabe}{Welcher Wert erfüllt $x \cdot (x + 7) = -12$? \\ \kreuz{$x = 3$} \kreuz{$x = 4$} \kreuz{$x = -4$} \kreuz{$x = -12$}}
\end{aufgabe}
\end{minipage}\par

% 11: Rechteckseite aus Fläche berechnen – flaechen-basis-k3-v1 (Original 2023-OS-B1c, ausgedacht: keine Marke)
\begin{aufgabe}{Ein rechteckiges Beet hat den Flächeninhalt $24$ m² und ist $4$ m breit. Wie lang ist es? \leerfeld[m]}
\end{aufgabe}

% 12: Term zu Sachtext angeben – terme-basis-k2-v7 (Original 2023-OS-B1h, ausgedacht: keine Marke)
\begin{aufgabe}{Die Summe aus einer Zahl $a$ und 6 wird verdoppelt. Welcher Term passt? \\ \kreuz{$2a + 6$} \kreuz{$2 \cdot (a + 6)$} \kreuz{$a + 6 \cdot 2$} \kreuz{$2 \cdot a \cdot 6$}}
\end{aufgabe}

% 13: Punktprobe durchführen – lineare-funktionen-basis-k3-v2 (Original 2023-OS-B1i, ausgedacht: keine Marke)
\begin{aufgabe}{Gegeben ist $f(x) = -3x + 4$. Welcher Punkt liegt nicht auf der Geraden? Kreuze an. \\ \kreuz{$(1|9)$} \kreuz{$(-2|10)$} \kreuz{$(3|-5)$}}
\end{aufgabe}

% 14: Eigenschaft einer Figur zuordnen – winkel-dreiecke-basis-k5-v2 (Original 2026-FOR-B1c, ausgedacht: keine Marke)
\begin{aufgabe}{Kreuze die richtige Ergänzung an: „In jedem Raute …“\\ \kreuz{… sind alle Winkel rechte Winkel.}\\ \kreuz{… sind alle Seiten gleich lang.}\\ \kreuz{… sind die Diagonalen gleich lang.}}
\end{aufgabe}

% 15: Arithmetisches Mittel berechnen – daten-basis-k1-v2 (Original 2021-OS-B1f, ausgedacht: keine Marke)
\begin{aufgabe}{In vier Spielen erzielt Ben $12$; $18$; $9$; $21$ Punkte. Berechne das arithmetische Mittel (Punktzahl). \leerfeld Punkte}
\end{aufgabe}

% 16: Bruchteil einer Fläche bestimmen – brueche-dezimalzahlen-basis-k1-v11 (Original 2026-FOR-B1b, ausgedacht: keine Marke)
\begin{aufgabe}{\begin{minipage}[t]{0.6\linewidth}Ein Fenster besteht aus gleich großen Scheiben; die grauen sind aus Milchglas. Wie viel Prozent der Fläche sind grau? \leerfeld[\%]\end{minipage}\hfill\begin{minipage}[t]{0.37\linewidth}\vspace{0pt}
\bruchrechteck[4]{6}{10}
\end{minipage}}
\end{aufgabe}

% 17: Gegenfläche im Würfelnetz bestimmen – koerper-basis-k1-v3 (Original 2016-OS-B1j, ausgedacht: keine Marke)
\begin{aufgabe}{\begin{minipage}[t]{0.6\linewidth}Das Netz aus sechs Quadraten (je Buchstabe ein Quadrat) wird zu einem Würfel gefaltet. Der Würfel wird auf die Fläche K gestellt. Welche Fläche liegt oben? \leerfeld\end{minipage}\hfill\begin{minipage}[t]{0.37\linewidth}\vspace{0pt}
$\begin{array}{cccc} K &  &  &  \\ L & M & N &  \\  &  & O & P \end{array}$
\end{minipage}}
\end{aufgabe}

% 18: Winkelfunktion Seitenverhältnis angeben – trigonometrie-basis-k2-v6 (Original 2025-OS-B1g, ausgedacht: keine Marke)
\begin{aufgabe}{\begin{minipage}[t]{0.6\linewidth}Rechter Winkel bei $C$. Trage den Bruch ein. $\mathrm{sin}\,\alpha =$ \leerfeld\end{minipage}\hfill\begin{minipage}[t]{0.37\linewidth}\vspace{0pt}
\dreieck{(0,0)}{(3.75,0)}{(1.35,1.8)}{x}{y}{z}{\alpha}{}{}
\end{minipage}}
\end{aufgabe}

\begleitteil
\blattkopf{Basisaufgaben · Zettel 13}{BAS-Z13 · Lösungen}
\erg{1}{$(-2{,}1 + (-2)) : 2 = -2{,}05$ \hfill {\footnotesize\textit{Mitte zweier Zahlen bestimmen · 2020}}}
\erg{2}{$-3 \cdot (-5 + 2) = -3 \cdot (-3) = 9$ \hfill {\footnotesize\textit{Termwert berechnen · 2026}}}
\erg{3}{z. B. $-1{,}5$; richtig ist jede Zahl größer als $-2$ und kleiner als $-1$ \hfill {\footnotesize\textit{Zahl zu Bedingung angeben · 2015}}}
\erg{4}{$-2$ – das Komma wandert zwei Stellen nach rechts \hfill {\footnotesize\textit{Zehnerpotenzschreibweise umwandeln · 2025}}}
\erg{5}{$36 : 400 = 0{,}09 = 9\,\%$ \hfill {\footnotesize\textit{Prozentsatz berechnen · 2014}}}
\erg{6}{$0{,}6 \cdot 25 = 15\,€$ \hfill {\footnotesize\textit{Prozentwert berechnen · 2026}}}
\erg{7}{$1{,}5$ kg $= 1500$ g, also $1{,}5$ kg $=$ $1500$ g \hfill {\footnotesize\textit{Größen vergleichen · 2026}}}
\erg{8}{$2$ l $= 2000$ ml; $2000 : 250 = 8$ \hfill {\footnotesize\textit{Portionen aus Gesamtmenge berechnen · 2022}}}
\erg{9}{Für $1$: $9 : 5 = 1{,}80$; $1{,}80 \cdot 8 = 14{,}40\,€$ \hfill {\footnotesize\textit{Proportionale Zuordnung Dreisatz · 2023}}}
\erg{10}{$x = -4$, denn $(-4) \cdot (-4 + 7) = (-4) \cdot 3 = -12$ \hfill {\footnotesize\textit{Lösung durch Einsetzen prüfen · 2025}}}
\erg{11}{$24 : 4 = 6$ m \hfill {\footnotesize\textit{Rechteckseite aus Fläche berechnen · 2023}}}
\erg{12}{$2 \cdot (a + 6)$ \hfill {\footnotesize\textit{Term zu Sachtext angeben · 2023}}}
\erg{13}{$(1|9)$, denn $f(1) = 1$ \hfill {\footnotesize\textit{Punktprobe durchführen · 2023}}}
\erg{14}{… sind alle Seiten gleich lang. \hfill {\footnotesize\textit{Eigenschaft einer Figur zuordnen · 2026}}}
\erg{15}{$60 : 4 = 15$ Punkte \hfill {\footnotesize\textit{Arithmetisches Mittel berechnen · 2021}}}
\erg{16}{$6$ von $10$ gleichen Teilen: $\frac{6}{10} = 60\,\%$ \hfill {\footnotesize\textit{Bruchteil einer Fläche bestimmen · 2026}}}
\erg{17}{O (O und K liegen sich gegenüber) \hfill {\footnotesize\textit{Gegenfläche im Würfelnetz bestimmen · 2016}}}
\erg{18}{$\mathrm{sin}\,\alpha = \frac{x}{z}$ \hfill {\footnotesize\textit{Winkelfunktion Seitenverhältnis angeben · 2025}}}
\end{document}
